% paper/sections/02_governing.tex
\section{支配方程式系と表面張力項}
\label{sec:governing}

\subsection{One-Fluid 定式化と非圧縮条件}

非圧縮条件は連続の式：
\begin{equation}
  \bnabla \cdot \bu = 0
  \label{eq:continuity}
\end{equation}
で与えられる．
物性値は界面指標 $\psi \in [0,1]$ により連続的に補間する：
\begin{align}
 \rho(\psi)&=\rho_g+(\rho_l-\rho_g)\,\psi, & \mu(\psi)&=\mu_g+(\mu_l-\mu_g)\,\psi.
\end{align}
ここで $\psi$ は Conservative Level Set 法におけるヘヴィサイド関数 $H_\varepsilon(\phi)$ に相当する．

\subsection{主要な変数と記法の定義}
\label{sec:notation}

本稿では，界面を表現するために二つの関連する変数，符号付き距離関数 $\phi$ と保存形レベルセット変数 $\psi$ を用いる．両者の役割を明確に区別することが，アルゴリズムを理解する上で極めて重要である．

\begin{tcolorbox}[defbox, title={変数 $\phi$ と $\psi$ の役割分担}]
\begin{itemize}
  \item \textbf{符号付き距離関数 $\phi(\bm{x}, t)$}:
        界面からの最短距離に符号を付けた関数（$\phi=0$ が界面）．
        その性質 $|\bnabla\phi|=1$ は，法線ベクトル $\hat{\bm{n}} = \bnabla\phi$ や曲率 $\kappa = -\bnabla\cdot\hat{\bm{n}}$ といった\textbf{幾何学的情報}の計算精度を保証するために不可欠である．
        本稿では，曲率計算や再初期化方程式における法線ベクトルの算出に $\phi$ を用いる．

  \item \textbf{保存形レベルセット変数 $\psi(\bm{x}, t)$}:
        $\phi$ を滑らかなヘヴィサイド関数 $H_\varepsilon$ で変換したもので，液相で 1，気相で 0 を取る．
        その移流方程式は保存形で書けるため，\textbf{質量（体積）保存性}を向上させる役割を担う．
        密度 $\rho(\psi)$ や粘性 $\mu(\psi)$ などの\textbf{物理的物性値}の補間にも直接用いられる．
\end{itemize}
両者は $\psi = H_\varepsilon(\phi)$ の関係で結ばれる．逆に $\psi$ から $\phi$ を得るには，非線形方程式 $H_\varepsilon(\phi) - \psi = 0$ をニュートン法などの反復解法で数値的に解く必要がある．
\end{tcolorbox}

\subsection{無次元化 NS 方程式}

参照量として液相密度 $\rho_l$，代表速度 $U$，代表長さ $L$ を選ぶ．
これらを用いて，物理変数を無次元化する．次元を持つ変数を $(\cdot)^*$，無次元変数を $(\cdot)$ と表すと，以下の関係式が成り立つ：

\begin{equation}
 \bm{x}^* = L\bm{x}, \quad \bu^* = U\bu, \quad t^* = \frac{L}{U}t, \quad p^* = \rho_l U^2 p
\end{equation}

これらを次元を持つ運動量保存則に代入し，整理することで，流体の挙動を支配する3つの主要な無次元パラメータが自然に導かれる．

\begin{itemize}
  \item \textbf{レイノルズ数 ($Re$)}：慣性力と粘性力の比
  \[ Re \equiv \frac{\rho_l U L}{\mu_l} \quad \left(\text{係数項：}\frac{\mu_l}{\rho_l U L} = \frac{1}{Re}\right) \]
  \item \textbf{フルード数 ($Fr$)}：慣性力と重力の比
  \[ Fr \equiv \frac{U}{\sqrt{gL}} \quad \left(\text{係数項：}\frac{gL}{U^2} = \frac{1}{Fr^2}\right) \]
  \item \textbf{ウェーバー数 ($We$)}：慣性力と表面張力の比
  \[ We \equiv \frac{\rho_l U^2 L}{\sigma} \quad \left(\text{係数項：}\frac{\sigma}{\rho_l U^2 L} = \frac{1}{We}\right) \]
\end{itemize}

具体的な導出過程は以下の通りである．次元系の運動量方程式の各項に変換式を代入する．
例えば，左辺の慣性項（非定常項）は次のように変換される：

\[ \rho^* \pd{\bu^*}{t^*} = (\rho_l\rho)\,\pd{(U\bu)}{\frac{L}{U}t} = \frac{\rho_l U^2}{L}\,\rho\,\pd{\bu}{t} \]

同様に，圧力勾配項は $(\rho_l U^2/L)\bnabla p$，粘性項は $(\mu_l U/L^2)\bnabla\cdot(\mu \dots)$ のように変換される．
方程式全体を慣性力の代表スケール $\rho_l U^2/L$ で除算すると，
\[
 \rho(\psi)\!\left(\pd{\bu}{t}+\bu\cdot\bnabla\bu\right)
 = -\bnabla p + \frac{\mu_l}{\rho_l U L}\bnabla\cdot\!\left[\mu(\psi)(\bnabla\bu+\bnabla\bu^T)\right]
 - \frac{gL}{U^2}\frac{\rho(\psi)}{\rho_l}\hat{\bm z} + \frac{\sigma}{\rho_l U^2 L}\kappa\,\bnabla\psi
\]
となり，前述の無次元パラメータが係数として現れる．Conservative Level Set (CLS) 法に基づく最終的な無次元形は以下となる：

\begin{tcolorbox}[mybox, title={無次元 Navier--Stokes（CLS 版）}]
\begin{equation}
 \rho(\psi)\!\left(\pd{\bu}{t}+\bu\cdot\bnabla\bu\right)
 =-\bnabla p+\frac{1}{Re}\,\bnabla\cdot\!\left[\mu\,(\bnabla\bu+\bnabla\bu^T)\right]
 -\frac{\rho}{Fr^2}\hat{\bm z}+\frac{\kappa\,\bnabla\psi}{We}
 \label{eq:NS_dimless_cls}
\end{equation}
\end{tcolorbox}

\subsubsection{CSF の正規化（Balanced-Force）}

Balanced-force の観点から，表面張力項 $\kappa\,\bnabla\psi/We$ は圧力勾配項と同一点・同重みで離散化する必要がある．
これにより，静止流体中での寄生流れ（parasitic currents）を抑制する．

\subsection{曲率の計算}
界面形状に関する幾何量（法線ベクトル・曲率）は，符号付き距離関数 $\phi$ から計算する方が精度が良い．
\begin{equation}
 \kappa(\phi) = -\bnabla\cdot \hat{\bm{n}} = -\bnabla\cdot\!\left(\frac{\bnabla\phi}{\norm{\bnabla\phi}}\right)
 \label{eq:curvature_div}
\end{equation}
この発散形式のまま離散化することも可能だが，微分誤差を最小化するために
商の微分公式を用いて展開した形を用いることが推奨される．

2次元直交座標系 $(x,y)$ において，法線ベクトルは $\hat{\bm{n}} = (\phi_x/|\bnabla\phi|, \phi_y/|\bnabla\phi|)$ と表される．
これを発散の定義式に代入して展開する：
\begin{align}
  \kappa &= -\left( \pd{}{x}\frac{\phi_x}{\sqrt{\phi_x^2+\phi_y^2}} + \pd{}{y}\frac{\phi_y}{\sqrt{\phi_x^2+\phi_y^2}} \right) \notag \\
  &= -\left[ \frac{\phi_{xx}(\phi_x^2+\phi_y^2)^{1/2} - \phi_x \cdot \frac{1}{2}(\phi_x^2+\phi_y^2)^{-1/2}(2\phi_x\phi_{xx} + 2\phi_y\phi_{xy})}{\phi_x^2+\phi_y^2} \right] - (\text{$y$成分}) \notag \\
  &= -\frac{\phi_{xx}(\phi_x^2+\phi_y^2) - \phi_x^2\phi_{xx} - \phi_x\phi_y\phi_{xy}}{(\phi_x^2+\phi_y^2)^{3/2}}
     -\frac{\phi_{yy}(\phi_x^2+\phi_y^2) - \phi_y^2\phi_{yy} - \phi_x\phi_y\phi_{xy}}{(\phi_x^2+\phi_y^2)^{3/2}} \notag \\
  &= -\frac{\phi_{yy}\phi_x^2 - 2\phi_x\phi_y\phi_{xy} + \phi_{xx}\phi_y^2}{(\phi_x^2+\phi_y^2)^{3/2}}
\end{align}
整理すると，以下の実装形式が得られる．

\begin{tcolorbox}[mybox, title={2次元界面曲率（展開形）}]
\begin{equation}
  \kappa = -\frac{
    \phi_y^2\,\phi_{xx}
    - 2\phi_x\phi_y\,\phi_{xy}
    + \phi_x^2\,\phi_{yy}
  }{\bigl(\phi_x^2 + \phi_y^2\bigr)^{3/2}}
  \label{eq:curvature_2d}
\end{equation}

\subsubsection{曲率計算の数値的安定化}

式 \eqref{eq:curvature_2d} は数学的には厳密だが，数値計算で直接使用するには注意が必要である．
分母の $\norm{\bnabla\phi}^{3/2}$ は，$\phi$ が完全な符号付き距離関数であれば界面近傍で 1 に近いため問題ない．
しかし，界面から離れた領域や，移流・再初期化の誤差によって $\norm{\bnabla\phi}$ がゼロに近づくと，
この項は発散し，計算が破綻する可能性がある．

\begin{tcolorbox}[warnbox, title={実装上の重要事項：ゼロ除算の回避}]
この問題を回避するため，分母の大きさに下限値を設けるのが一般的である．
具体的には，小さな正の定数 $\epsilon_{\text{norm}}$（例：$10^{-8}$）を導入し，
分母を $\max(\norm{\bnabla\phi}, \epsilon_{\text{norm}})^{3}$ で置き換える．
\[ \kappa \approx -\frac{ \phi_y^2\,\phi_{xx} - 2\phi_x\phi_y\,\phi_{xy} + \phi_x^2\,\phi_{yy} }{ \max\bigl(\phi_x^2 + \phi_y^2, \epsilon_{\text{norm}}^2\bigr)^{3/2} } \]
この処理により，$\bnabla\phi$ がゼロになる点（特異点）での発散を防ぎ，計算のロバスト性を確保できる．
さらに，計算された $\kappa$ に対して微小なガウシアンフィルタを適用し，格子スケールの数値ノイズを除去することも有効な手段である．
\end{tcolorbox}
\end{tcolorbox}
